\chapter{Commutative algebra over valuation and graded rings}
\index{valuation ring!commutative algebra}\index{graded ring}
\label{app:valuation-algebra}

This appendix collects the algebra used to pass between ideals, valuations,
blowups, and finite morphisms.  It supplies proofs of elementary facts often
compressed into an appeal to the Rees algebra, and separates three notions
that must not be conflated: integral closure, subintegrality, and universal
homeomorphy.

Unless otherwise stated, rings are commutative with unit and ideals are
proper.  A domain has fraction field \(K\).  The noetherian hypothesis is
always displayed when it is used.  Finiteness of normalization is invoked
only for Nagata rings, in particular for finitely generated algebras over a
field.

\section{Valuation rings and ideal values}
\label{sec:appA-valuations}

\begin{definition}[Krull valuation]
\label{def:appA-krull-valuation}
Let \(\Gamma\) be a totally ordered abelian group and
\(\Gamma_\infty=\Gamma\cup\{+\infty\}\).
A valuation on a field \(K\) is a map
\[
 v:K\longrightarrow\Gamma_\infty
\]
such that
\[
 v(0)=+\infty,\qquad v(K^\times)\subset\Gamma,\qquad v(1)=0,
\]
\[
 v(fg)=v(f)+v(g),\qquad
 v(f+g)\geq\min\{v(f),v(g)\}.
\]
Its valuation ring and maximal ideal are
\[
 R_v=\{f\in K:v(f)\geq0\},\qquad
 \mathfrak m_v=\{f\in K:v(f)>0\}.
\]
\end{definition}

\begin{lemma}\label{lem:appA-valuation-dichotomy}
A subring \(R\subset K\) is a valuation ring of \(K\) if and only if, for
each \(f\in K^\times\), either \(f\in R\) or \(f^{-1}\in R\).
\end{lemma}

\begin{proof}
For \(R=R_v\), either \(v(f)\geq0\) or \(v(f^{-1})=-v(f)>0\).
Conversely, the principal fractional ideals of a ring with the stated
dichotomy are totally ordered: if \(a,b\ne0\), then either \(a/b\in R\) or
\(b/a\in R\).  The ordered group \(K^\times/R^\times\), with order induced
by reverse inclusion of principal fractional ideals, yields the valuation.
\end{proof}

Let \(A\subset K\), and suppose \(v\) is nonnegative on \(A\).  For a
nonzero finitely generated ideal
\(\mathfrak a=(f_1,\ldots,f_r)\subset A\), put
\begin{equation*}
 v(\mathfrak a):=\min_i v(f_i).
 \tag{A.1}\label{eq:appA-ideal-value}
\end{equation*}
This is independent of the generators.  Every element
\(\sum a_if_i\) has value at least \(\min_i v(f_i)\), and each \(f_i\)
belongs to the ideal.

\begin{proposition}\label{prop:appA-ideal-arithmetic}
For nonzero finitely generated ideals \(\mathfrak a,\mathfrak b\subset A\),
\[
 v(\mathfrak a\mathfrak b)=v(\mathfrak a)+v(\mathfrak b),\qquad
 v(\mathfrak a+\mathfrak b)
   =\min\{v(\mathfrak a),v(\mathfrak b)\}.
\]
If \(\mathfrak a\subset\mathfrak b\), then
\(v(\mathfrak a)\geq v(\mathfrak b)\).
\end{proposition}

\begin{proof}
Generators of the product are the pairwise products of generators, so the
first identity follows from additivity.  The union of two generating sets
generates the sum, proving the second identity.  The last assertion follows
directly from \eqref{eq:appA-ideal-value}.
\end{proof}

\begin{remark}[Checkpoint]
\label{rem:appA-ideal-sign}
The book uses \(\phi_{\mathfrak a}(v)=-v(\mathfrak a)\).  Thus ideal
inclusion and function order point in opposite directions:
\[
 \mathfrak a\subset\mathfrak b
 \quad\Longrightarrow\quad
 \phi_{\mathfrak a}\leq\phi_{\mathfrak b}.
\]
This sign is first used for flag ideals in
\Cref{chap:bj-test-configurations} and is essential in
\Cref{chap:bj-smooth-characteristic}.
\end{remark}

\section{Integral dependence and integral closure}
\label{sec:appA-integral-dependence}

\begin{definition}[Integral element and extension]
\label{def:appA-integral-element}
An element \(b\) of an \(A\)-algebra \(B\) is integral over \(A\) if
\[
 b^d+a_1b^{d-1}+\cdots+a_d=0,\qquad a_i\in A.
\]
The extension \(A\to B\) is integral if every element of \(B\) is integral.
The integral closure of a domain \(A\) in an algebraic extension \(L/K\)
is the set of elements of \(L\) integral over \(A\).
\end{definition}

\begin{proposition}\label{prop:appA-finite-module-criterion}
For \(b\in B\), the following are equivalent:
\begin{enumerate}
\item \(b\) is integral over \(A\);
\item \(A[b]\) is a finite \(A\)-module;
\item there exists a faithful finite \(A[b]\)-module that is finite as an
      \(A\)-module.
\end{enumerate}
\par
\end{proposition}

\begin{proof}
A monic equation expresses every \(b^m\), \(m\geq d\), as an
\(A\)-linear combination of \(1,b,\ldots,b^{d-1}\), proving
\((1)\Rightarrow(2)\).  The implication \((2)\Rightarrow(3)\) uses
\(A[b]\).  For \((3)\Rightarrow(1)\), multiplication by \(b\) on a finite
set of \(A\)-generators and the determinant trick yield a monic polynomial
that kills the faithful module, hence kills \(b\).
\end{proof}

\begin{definition}[Integral closure of an ideal]
\label{def:appA-integral-closure-ideal}
Let \(\mathfrak a\subset A\).  An element \(f\in A\) is integral over
\(\mathfrak a\) if
\[
 f^d+c_1f^{d-1}+\cdots+c_d=0,
 \qquad c_i\in\mathfrak a^i.
\]
The set of such elements is the integral closure
\(\overline{\mathfrak a}\).
\end{definition}

\begin{lemma}\label{lem:appA-rees-translation}
For an indeterminate \(T\), an element \(f\in A\) is integral over
\(\mathfrak a\) if and only if \(fT\) is integral over
\[
 \mathcal R_A(\mathfrak a)
   :=A[\mathfrak aT]
   =\bigoplus_{m\geq0}\mathfrak a^mT^m
   \subset A[T].
\]
\end{lemma}

\begin{proof}
Multiply an integral equation for \(f\) by \(T^d\).  Conversely, take the
degree-\(d\) homogeneous part of a monic equation for \(fT\), and set
\(T=1\).
\end{proof}

\begin{theorem}\label{thm:appA-valuative-integral-closure}
Let \(A\) be a noetherian domain with fraction field \(K\), and let
\(\mathfrak a\ne0\).  Then
\[
 \overline{\mathfrak a}
   =\{f\in A:v(f)\geq v(\mathfrak a)
       \text{ for every valuation }v\text{ of }K
       \text{ nonnegative on }A\}.
\tag{A.2}\label{eq:appA-valuative-integral-closure}
\]
It is enough to test the finitely many discrete rank-one Rees valuations of
\(\mathfrak a\).
\end{theorem}

\begin{proof}[From the valuative integral-closure input]
If \(f\) satisfies an integral equation and
\(v(f)<v(\mathfrak a)\), then
\[
 v(f^d)<v(c_if^{d-i})
\]
for every \(i\geq1\).  A sum cannot vanish when exactly one summand has
strictly smallest value, a contradiction.  The converse and finite
reduction to Rees valuations are the foundational valuative theorem for
integral dependence; we import precisely this converse and finite-reduction
statement from \cite[Theorem~6.8.3 and Theorem~10.2.2]{SwansonHuneke2006};
see also \cite[Chapter~VI]{ZariskiSamuelII}.  The Rees valuations
are the valuations of the prime components of the Cartier divisor cut out
by \(\mathfrak a\) on the normalized blowup.  Components dominating a
divisor of \(\Spec A\) need not be exceptional.
\end{proof}

\begin{corollary}\label{cor:appA-integral-closure-values}
For every valuation centered on \(A\),
\[
 v(\overline{\mathfrak a})=v(\mathfrak a).
\]
\end{corollary}

\begin{proof}
The inclusion \(\mathfrak a\subset\overline{\mathfrak a}\) gives one
inequality; \Cref{thm:appA-valuative-integral-closure} gives the
other for each element of \(\overline{\mathfrak a}\).
\end{proof}

\begin{example}[A monomial computation]
\label{ex:appA-monomial-integral-closure}
In \(k[x,y]\), set \(\mathfrak a=(x^2,y^3)\).
For the monomial valuation \(v_{\alpha,\beta}\),
\[
 v_{\alpha,\beta}(x^iy^j)=i\alpha+j\beta.
\]
The inequalities against \(v_{\alpha,\beta}(\mathfrak a)\), for all
\(\alpha,\beta>0\), are equivalent to \(i/2+j/3\geq1\).  Thus
\[
 \overline{(x^2,y^3)}=(x^2,xy^2,y^3).
\]
Here is a direct check that avoids assuming a monomial theorem.  The
monomial valuation with \((\alpha,\beta)=(3,2)\) assigns value \(6\) to
\(\mathfrak a\).  For a polynomial it takes the minimum of \(3i+2j\)
over its nonzero monomials, so every element of
\(\overline{\mathfrak a}\) consists of monomials satisfying
\(3i+2j\geq6\).  For nonnegative integer exponents, these monomials
are exactly those divisible by \(x^2\), \(xy^2\), or \(y^3\).
Conversely, the two original generators are integral, and \(z=xy^2\)
satisfies \(z^2-x^2y^4=0\), with
\(x^2y^4=y(x^2y^3)\in\mathfrak a^2\).  Thus \(xy^2\) is integral
over \(\mathfrak a\), proving both inclusions.  The Newton polygon is
the convex shadow of \eqref{eq:appA-valuative-integral-closure}.
\end{example}

\section{Rees algebras, filtrations, and blowups}
\label{sec:appA-rees-algebras}

A multiplicative filtration \(F^\lambda R_m\) has a Rees algebra even
without finite generation.  A finite-type test configuration requires that
Rees algebra to be finitely generated.  For rational weights with a common
denominator, clearing that denominator gives the bigraded algebra
\[
 \mathcal R(F)=
 \bigoplus_{m\geq0}\ \bigoplus_{\ell\in\mathbf Z}
 F^\ell R_m\,T^{-\ell}.
\tag{A.3}\label{eq:appA-filtered-rees}
\]
The gradings record section degree and weight.

\begin{proposition}\label{prop:appA-blowup-proj}
For a finite-type ideal \(\mathcal I\) on \(X\),
\[
 \operatorname{Bl}_{\mathcal I}X
   =\operatorname{Proj}_X
       \bigoplus_{m\geq0}\mathcal I^m.
\tag{A.4}\label{eq:appA-blowup-proj}
\]
The pullback ideal is invertible, and the blowup is universal among
morphisms \(g:Y\to X\) for which \(\mathcal I\cO_Y\) is invertible.
\end{proposition}

\begin{proof}
On \(X=\Spec A\), the standard chart \(D_+(fT)\), \(f\in\mathfrak a\), is
\(\Spec A[\mathfrak a/f]\), and the extended ideal is generated by \(f\).
These charts glue.  Any morphism on which the ideal is locally principal
factors uniquely through the corresponding charts.
\end{proof}

\begin{definition}[Normalized blowup]
\label{def:appA-normalized-blowup}
If \(X\) is integral and \(\mathcal I\) is a nonzero finite-type ideal,
the normalized blowup is the normalization, in the
unchanged function field \(K(X)\), of
\(\operatorname{Bl}_{\mathcal I}X\).  If \(X\) is Nagata, normalization is
finite by
\cite[\href{https://stacks.math.columbia.edu/tag/035S}{Tag 035S}]
{stacks-project}.
\end{definition}

Suppose in addition that \(A\) is normal and \(\mathfrak a\ne0\).
The normalization of the Rees
algebra in \(K(T)\) is then the graded algebra
\[
 \bigoplus_{m\geq0}\overline{\mathfrak a^m}T^m
 \subset A[T].
\]
Hence, on a normal integral scheme,
\[
 \operatorname{Bl}_{\mathfrak a}^{\nu}X
 =\operatorname{Proj}_X
   \bigoplus_{m\geq0}\overline{\mathfrak a^m}.
\tag{A.5}\label{eq:appA-normalized-rees}
\]
This is the algebra behind normalized flag ideals in
\Cref{chap:bj-test-configurations}.

If \(A\) is not normal, the degree-zero part of the normalization is
\(A^\nu\), not \(A\).  Thus \eqref{eq:appA-normalized-rees} is not asserted
on a nonnormal affine chart.  All later uses of this formula, in particular
\Cref{thm:appB-rees-valuations}, are over a normal base.

\begin{remark}[Checkpoint]
\label{rem:appA-normalization-versus-saturation}
Replacing each homogeneous piece by its integral closure is compatible
with normalization in \eqref{eq:appA-normalized-rees}.  Replacing a Rees
algebra by an arbitrary saturation is a different operation.  In
particular, finite generation must be checked before a filtration is
declared to be a test configuration.
\end{remark}

\section{Normalization and the conductor}
\label{sec:appA-normalization-conductor}

Let \(A\) be reduced with total quotient ring \(Q(A)\), and let \(A^\nu\)
be its normalization.  The conductor is
\[
 \mathfrak c_{A^\nu/A}
 :=\{a\in A:aA^\nu\subset A\}
 =\operatorname{Ann}_A(A^\nu/A).
\tag{A.6}\label{eq:appA-conductor}
\]
It is an ideal of both \(A\) and \(A^\nu\).

\begin{lemma}\label{lem:appA-conductor-square}
If \(A^\nu\) is finite over \(A\), and
\(\mathfrak c=\mathfrak c_{A^\nu/A}\), then
\[
\begin{tikzcd}
A \arrow[r] \arrow[d] &
A^\nu \arrow[d]\\
A/\mathfrak c \arrow[r] &
A^\nu/\mathfrak c
\end{tikzcd}
\]
is cartesian.
\end{lemma}

\begin{proof}
If \(b\in A^\nu\) has the same residue as \(a\in A\), then
\(b-a\in\mathfrak c\subset A\), hence \(b\in A\).  The other inclusion is
immediate.
\end{proof}

Away from \(V(\mathfrak c)\), normalization is an isomorphism.  The
conductor does not imply bijectivity over \(V(\mathfrak c)\); that is a
separate unibranch or subintegrality issue.

\section{Subintegrality and seminormality}
\label{sec:appA-subintegrality}

\begin{definition}[Subintegral extension]
\label{def:appA-subintegral}
An integral extension \(A\subset B\) is subintegral if
\(\Spec B\to\Spec A\) is bijective and induces an isomorphism of residue
fields at corresponding primes.
\end{definition}

\begin{definition}[Seminormality]
\label{def:appA-seminormal}
A reduced ring \(A\) is seminormal in an integral overring \(B\) if
\[
 b\in B,\quad b^2,b^3\in A
 \quad\Longrightarrow\quad b\in A.
\tag{A.7}\label{eq:appA-square-cube}
\]
It is seminormal if it is seminormal in its normalization.
\end{definition}

\begin{theorem}\label{thm:appA-square-cube}
For an integral extension \(A\subset B\), subintegrality is equivalent to
the following finite-stage condition: every finite subset of \(B\) lies in
a tower of elementary extensions \(C\subset C[b]\) with
\(b^2,b^3\in C\).  This construction is stable under polynomial base
change.
\end{theorem}

\begin{proof}[Foundational input]
We import exactly this finite-stage characterization from
\cite[\href{https://stacks.math.columbia.edu/tag/0CND}{Tag 0CND}]
{stacks-project}, with base-change statement
\cite[\href{https://stacks.math.columbia.edu/tag/0EUH}{Tag 0EUH}]
{stacks-project}.  The book uses it for finite subintegral extensions in
\Cref{chap:bj-subintegral-normalization}.
\end{proof}

\begin{example}[The cusp]
\label{ex:appA-cusp-seminormalization}
For \(A=k[t^2,t^3]\subset B=k[t]\), one has \(t^2,t^3\in A\) but
\(t\notin A\).  Thus \(A\) is not seminormal.  The normalization map is
subintegral: it changes the structure sheaf but neither points nor residue
fields.
\end{example}

\begin{definition}[Universal homeomorphism]
\label{def:appA-universal-homeomorphism}
A morphism is radicial if it is universally injective.  It is a universal
homeomorphism if every base change induces a homeomorphism of spaces.
\end{definition}

\begin{theorem}\label{thm:appA-universal-homeomorphism-criterion}
A morphism is a universal homeomorphism if and only if it is integral,
surjective, and radicial.
\end{theorem}

\begin{proof}[Foundational input]
This equivalence is exactly
\cite[\href{https://stacks.math.columbia.edu/tag/04DF}{Tag 04DF}]
{stacks-project}.
\end{proof}

\begin{remark}[Subintegral is stronger]
\label{rem:appA-subintegral-stronger}
A subintegral extension yields a universal homeomorphism.  The converse
fails: Frobenius in characteristic \(p\) is a universal homeomorphism but
need not induce identity residue extensions.  In characteristic zero, a
finite universal homeomorphism has trivial residue extensions, and it is
subintegral as soon as the ring map \(A\to B\) is injective --- for
instance whenever \(A\) is reduced, since the kernel lies inside every
prime of \(A\), hence inside its nilradical.  Injectivity cannot be dropped,
\Cref{def:appA-subintegral}
being a condition on an extension: \(\mathbf Q[x]/(x^2)\to\mathbf Q\) is a
finite universal homeomorphism with trivial residue extension and nonzero
kernel.
\end{remark}

\section{Determinants and norm polynomial laws}
\label{sec:appA-determinant-norms}

Let \(A\to B\) be finite locally free of rank \(d\).  For \(b\in B\), put
\[
 \operatorname N_{B/A}(b):=\det(m_b),
\tag{A.8}\label{eq:appA-element-norm}
\]
where \(m_b\) is multiplication by \(b\).

\begin{proposition}\label{prop:appA-norm-identities}
For \(b,c\in B\),
\[
 \operatorname N(1)=1,\qquad
 \operatorname N(bc)=\operatorname N(b)\operatorname N(c).
\]
The norm commutes with arbitrary base change.  For indeterminates
\(\mathbf T\), the expression
\[
 \operatorname N\left(\sum_jT_jb_j\right)
\]
is homogeneous of degree \(d\) in \(\mathbf T\).
\end{proposition}

\begin{proof}
Use \(m_1=I\), \(m_{bc}=m_bm_c\), and base change for determinants.
Every matrix entry of \(m_{\sum T_jb_j}\) is linear in the \(T_j\), so the
determinant has degree \(d\).
\end{proof}

For finite locally free \(q:W\to Z\) and a line bundle \(M\) on \(W\),
\[
 \operatorname N_q(M)
 =\det(q_*M)\otimes\det(q_*\cO_W)^{-1}.
\tag{A.9}\label{eq:appA-norm-line}
\]
Locally on \(Z\), trivialize \(M\) over the finite preimage.  Changing
a generator by \(u\in\cO_W^\times\) changes its determinant transition
by \(\operatorname N_q(u)\); the denominator cancels the underlying
finite-algebra determinant.  The element-norm identities therefore make
this line-bundle norm multiplicative and base-change compatible; see
\cite{FerrandNorm} and
\cite[\href{https://stacks.math.columbia.edu/tag/0BD2}{Tag 0BD2}]
{stacks-project}.  \Cref{chap:bj-finite-descent} uses the universal
polynomial in the \(T_j\), not a product over geometric points, because
fibres may be inseparable or nonreduced.

\begin{remark}[Checkpoint: flatten before taking norms]
\label{rem:appA-flatten-before-norms}
For a merely finite morphism, \(q_*\cO_W\) need not be locally free.
\Cref{chap:bj-finite-descent} first flattens the finite algebra and
only then takes determinant norms.  This order is part of the proof.
\end{remark}

\section{The Roberts--Singh input}
\label{sec:appA-roberts-singh}

\begin{theorem}\label{thm:appA-roberts-singh}
Let \(A\subset B\) be a subintegral extension of \(\mathbf Q\)-algebras
and let \(b\in B\).  There exist integers \(N\geq p\geq0\) and
\(c_1,\ldots,c_p\in B\), independent of \(m\), such that
\begin{equation*}
 b^m+\sum_{i=1}^{p}\binom mi c_i b^{m-i}\in A
 \qquad(m\geq N).
 \tag{A.10}\label{eq:appA-RS}
\end{equation*}
\end{theorem}

\begin{proof}[Foundational input and exact provenance]
\cite[Theorem~4.17]{RobertsSingh} makes every element subintegral in the
elementwise sense used there.
\cite[Proposition~4.2 and Corollary~4.3]{RobertsSingh} identify this with
the exponent-zero formula \eqref{eq:appA-RS}.  Their general systems
also have coefficients fixed in \(m\); exponent zero means that the
correction powers are \(b^{m-i}\), without the additional shift
\(b^{m-s-i}\) of a system of exponent \(s\).  Enlarging \(N\) to at least
\(p\) keeps every power \(b^{m-i}\) nonnegative, even when \(b=0\).
\end{proof}

\begin{corollary}\label{cor:appA-RS-generic}
For a finite subintegral extension of \(\mathbf Q\)-algebras and
\(g_j\in B[u]\), formula \eqref{eq:appA-RS} applies over
\(A[u,\mathbf z]\) to \(b=\sum_jz_jg_j\).
\end{corollary}

\begin{proof}
Polynomial base change preserves subintegrality by
\Cref{thm:appA-square-cube}; apply
\Cref{thm:appA-roberts-singh}.
\end{proof}

Taking homogeneous \(\mathbf z\)-coefficients and finite differences in
\(m\) produces the sandwich
\[
 (g_1,\ldots,g_s)^{n+2p}
 \subset K_nB[u]\subset(g_1,\ldots,g_s)^n.
\tag{A.11}\label{eq:appA-RS-sandwich}
\]
The calculation, including its two gradings and equivariance, is given in
\Cref{chap:bj-subintegral-normalization}.  Nothing here asserts
\eqref{eq:appA-RS} for arbitrary integral extensions.

\section{Dependency map}
\label{sec:appA-dependency-map}

\begin{center}
\begin{tabular}{p{0.33\textwidth}p{0.52\textwidth}}
\toprule
Tool & First substantial use\\
\midrule
Ideal values and integral closure
 & \Cref{chap:bj-test-configurations}\\
Normalized Rees algebra
 & \Cref{chap:bj-cartier-realization}\\
Finite normalization and conductor
 & \Cref{chap:bj-finite-descent,chap:bj-subintegral-normalization}\\
Universal-homeomorphism criterion
 & \Cref{chap:bj-universal-homeomorphisms}\\
Determinant norm
 & \Cref{chap:bj-finite-descent}\\
Roberts--Singh additive powers
 & \Cref{chap:bj-subintegral-normalization}\\
\bottomrule
\end{tabular}
\end{center}

\section{Source note}
\label{sec:appA-source-note}
Integral dependence and its valuative criterion follow
\cite[Chapters~V--VI]{ZariskiSamuelII}; noetherian normalization facts may
also be found in \cite{Mat89}.  Stable tags are audited in
\Cref{app:stacks}.  Norm functors are from \cite{FerrandNorm}.
The additive-power theorem is quoted only in the form proved in
\cite{RobertsSingh}.
